\chapter{Seguridad}

\marginal{el perímetro, enumerado}
Este capítulo usa un perímetro propio y lo enumera para que el lector pueda
rehacerlo: quince subfunciones de cinco funciones de la finalidad Gobierno, a
saber Justicia (subfunciones 1 a 4), Coordinación de la Política de Gobierno (2 a
5), Seguridad Nacional (1 a 3), Asuntos de Orden Público y de Seguridad Interior
(1 a 4) y Otros Servicios Generales (4 y 5).

Con ese perímetro, el gasto en seguridad pasa de \textbf{387,626.5 a 333,091.1
millones de pesos}: una caída real de \textbf{17.6\%}, la más grande de cualquier
rubro de este documento.

\small
\begin{longtable}{lSSSSS}
\caption{Perímetro de seguridad, quince subfunciones (mdp)}\label{tab:C10_1}\\
\toprule
{Concepto} & {2024} & {2025} & {Dif. nominal} & {Dif. real} & {Var. real \%} \\
\midrule\endfirsthead
\toprule
{Concepto} & {2024} & {2025} & {Dif. nominal} & {Dif. real} & {Var. real \%} \\
\midrule\endhead
\bottomrule\endfoot
{TOTAL perimetro de seguridad (15 subfunciones)} & \num{387626.5} & \num{333091.1} & \num{-54535.4} & \num{-71203.3} & \num{-17.6} \\
\multicolumn{6}{l}{\itshape por funcion y subfuncion} \\
{1.6.1 Seguridad Nacional} & \num{103530.1} & \num{101948.0} & \num{-1582.2} & \num{-6034.0} & \num{-5.6} \\
{1.2.1 Justicia} & \num{84778.9} & \num{77263.3} & \num{-7515.6} & \num{-11161.1} & \num{-12.6} \\
{1.7.1 Asuntos de Orden Público y de Segu} & \num{72976.2} & \num{35951.5} & \num{-37024.7} & \num{-40162.7} & \num{-52.8} \\
{1.6.2 Seguridad Nacional} & \num{44088.4} & \num{34303.9} & \num{-9784.4} & \num{-11680.2} & \num{-25.4} \\
{1.2.3 Justicia} & \num{25032.2} & \num{25242.8} & \num{210.6} & \num{-865.8} & \num{-3.3} \\
{1.2.2 Justicia} & \num{20527.8} & \num{21241.8} & \num{714.0} & \num{-168.7} & \num{-0.8} \\
{1.7.4 Asuntos de Orden Público y de Segu} & \num{11033.4} & \num{11396.6} & \num{363.2} & \num{-111.2} & \num{-1.0} \\
{1.7.3 Asuntos de Orden Público y de Segu} & \num{4000.0} & \num{5578.4} & \num{1578.4} & \num{1406.4} & \num{33.7} \\
{1.2.4 Justicia} & \num{5176.8} & \num{5046.4} & \num{-130.4} & \num{-353.0} & \num{-6.5} \\
{1.3.2 Coordinación de la Política de Gob} & \num{4390.7} & \num{4230.5} & \num{-160.2} & \num{-349.0} & \num{-7.6} \\
{1.8.5 Otros Servicios Generales} & \num{4320.4} & \num{3984.1} & \num{-336.4} & \num{-522.1} & \num{-11.6} \\
{1.3.4 Coordinación de la Política de Gob} & \num{3230.2} & \num{3235.5} & \num{5.4} & \num{-133.5} & \num{-4.0} \\
{1.6.3 Seguridad Nacional} & \num{2984.3} & \num{2760.0} & \num{-224.3} & \num{-352.6} & \num{-11.3} \\
{1.8.4 Otros Servicios Generales} & \num{1073.4} & \num{466.9} & \num{-606.5} & \num{-652.7} & \num{-58.3} \\
{1.7.2 Asuntos de Orden Público y de Segu} & \num{190.7} & \num{176.0} & \num{-14.6} & \num{-22.8} & \num{-11.5} \\
{1.3.5 Coordinación de la Política de Gob} & \num{164.5} & \num{149.6} & \num{-14.9} & \num{-21.9} & \num{-12.8} \\
{1.3.3 Coordinación de la Política de Gob} & \num{128.7} & \num{115.8} & \num{-12.9} & \num{-18.4} & \num{-13.7} \\
\multicolumn{6}{l}{\itshape por ramo} \\
{07 Defensa Nacional} & \num{104840.6} & \num{103213.6} & \num{-1627.0} & \num{-6135.1} & \num{-5.6} \\
{03 Poder Judicial} & \num{78327.3} & \num{70983.6} & \num{-7343.7} & \num{-10711.7} & \num{-13.1} \\
{36 Seguridad y Protección Ciudadana} & \num{105838.8} & \num{70422.2} & \num{-35416.6} & \num{-39967.7} & \num{-36.2} \\
{13 Marina} & \num{44101.4} & \num{34325.2} & \num{-9776.2} & \num{-11672.6} & \num{-25.4} \\
{49 Fiscalía General de la República} & \num{19333.4} & \num{20126.0} & \num{792.6} & \num{-38.7} & \num{-0.2} \\
{33 Aportaciones Federales para Entidades Federativas y Municipios} & \num{9210.9} & \num{9565.5} & \num{354.6} & \num{-41.4} & \num{-0.4} \\
{04 Gobernación} & \num{7919.4} & \num{6474.0} & \num{-1445.4} & \num{-1785.9} & \num{-21.6} \\
{06 Hacienda y Crédito Público} & \num{4694.1} & \num{4327.6} & \num{-366.5} & \num{-568.3} & \num{-11.6} \\
{32 Tribunal Federal de Justicia Administrativa} & \num{3304.5} & \num{3304.5} & \num{0.0} & \num{-142.1} & \num{-4.1} \\
{35 Comisión Nacional de los Derechos Humanos} & \num{1722.1} & \num{1722.4} & \num{0.2} & \num{-73.8} & \num{-4.1} \\
\end{longtable}
\par\vspace{2pt}\footnotesize Fuente: elaboración propia del ITED con los analíticos del PEF aprobado 2024 y 2025. Las diferencias en millones de pesos se reportan dos veces, nominal y real; la real está en pesos de 2025 con el deflactor de 4.3\% que declara el CGPE.
\normalsize


\section{Aquí sí es una decisión de gasto}

A diferencia de salud, esta caída no se disuelve al descomponerla.
\textbf{El movimiento está concentrado en una unidad responsable que no cambió de
ramo ni de nombre.} La Guardia Nacional, unidad H00 del Ramo 36, pasa de
\textbf{70,767.4 a 33,799.8 millones de pesos}: menos de la mitad.

Se comprobó que no se trata de una reubicación. En 2025 la Guardia Nacional sigue
siendo la unidad H00 del Ramo 36 y \textbf{no aparece ninguna unidad con ese
nombre en la Secretaría de la Defensa Nacional}. La subfunción de asuntos de orden
público que la contiene cae 52.8\%, de 72,976.2 a 35,951.5.

Los demás movimientos del Ramo 36 apuntan en la misma dirección: las provisiones
para infraestructura de seguridad caen de 23,272.7 a 10,208.9, y el apoyo
administrativo de 13,889.1 a 9,284.7.

\section{Lo militar no cae}

Conviene decirlo porque el capítulo 4 mostró que el Ramo 07 pierde 112,301.7
millones de pesos en términos reales. \textbf{Ese movimiento no es seguridad}: es
la obra ferroviaria y aeroportuaria que Defensa deja de ejecutar. La función
Seguridad Nacional, que es donde vive el gasto militar propiamente dicho, cae
5.6\% real en su primera subfunción y 25.4 en la segunda, muy por debajo de la
caída del perímetro completo.

\textbf{La lectura del paquete 2025 en seguridad es que se reduce el componente
civil, y en particular el cuerpo policial federal, mientras el componente militar
se mantiene.} Es una afirmación fuerte y por eso este capítulo enumeró su
perímetro y verificó la ubicación de la unidad responsable antes de hacerla.

\section{Lo que falta}

El paquete no publica estado de fuerza, despliegue ni indicadores de resultado, de
modo que este capítulo puede decir cuánto y quién, y no puede decir con cuántos
elementos ni con qué efecto. La pregunta sobre la obra pública no militar
ejecutada por las fuerzas armadas se responde en el capítulo 8.

\clearpage
